\chapter{Nilai Eigen dan Diagonalisasi}

Setiap menit, sebuah pemantau jaringan membagi seluruh permintaan yang masuk ke dalam tiga
keranjang: host A, host B, dan host C. Dari catatan lama ia menyusun aturan perpindahan. Dari
permintaan yang menuju A, 80 persen tetap di A, 10 persen pindah ke B, dan 10 persen pindah ke C.
Dari B, 30 persen pindah ke A, 60 persen bertahan, dan 10 persen pindah ke C. Dari C, 20 persen
pindah ke A, 30 persen ke B, dan 50 persen bertahan. Bila bagian permintaan di ketiga keranjang
pada menit ke-$k$ ditulis sebagai kolom $\mathbf{x}_{k}$, aturan itu berbentuk
$\mathbf{x}_{k+1}=M\mathbf{x}_{k}$ dengan
\[
  M=\begin{pmatrix}0.8&0.3&0.2\\ 0.1&0.6&0.3\\ 0.1&0.1&0.5\end{pmatrix}.
\]

Mulai dari $\mathbf{x}_{0}=(1,0,0)$, yaitu semua permintaan mula-mula menuju A, pemantau itu
menghitung sepuluh langkah dan memperoleh
$\mathbf{x}_{10}=(0.5672,\;0.2661,\;0.1666)$. Tiga puluh langkah kemudian angkanya
$(0.5667,\;0.2667,\;0.1667)$. Jumlah ketiga entri tetap satu, dan pada pembacaan kedua angka itu
berhenti bergerak sampai empat angka di belakang koma.

Pertanyaan yang layak diajukan bukan berapa hasilnya setelah empat puluh menit; perkalian berulang
sudah menjawabnya untuk satu titik awal. Pertanyaan yang lebih berguna: adakah vektor yang hanya
diperbesar oleh $M$ tanpa diubah arahnya? Di sini arah itu adalah
$(17/30,\;4/15,\;1/6)=(0.5667,\;0.2667,\;0.1667)$, tepat angka yang muncul setelah empat puluh
langkah, dan pengalinya $1$.

Bab 8 memberi determinan: $\det M=1/5$, jadi $M$ invertible dan setiap sistem $M\mathbf{y}=
\mathbf{b}$ punya solusi tunggal. Itu menjawab apakah proses dapat dibalik, bukan ke arah mana ia
menuju, dan bukan seberapa cepat. Angka yang mengatur laju di sini adalah $1/2$: nilai eigen
terbesar setelah $1$. Karena $(1/2)^{40}\approx9.1\times10^{-13}$, galat awal tinggal di belakang
enam angka pertama, dan itulah sebabnya pembacaan kedua terlihat diam.

Bab ini menyiapkan alat untuk membaca dua angka itu langsung dari matriksnya: cara menghitungnya,
syarat agar arah yang dipertahankan cukup banyak untuk menjadi basis, apa yang berubah ketika
basis diganti, lalu tiga pemakaiannya: pangkat matriks, sistem diferensial linear, dan rantai
Markov.

\begin{untukanak}
Ambil selembar karet yang diberi garis-garis kotak, lalu tarik dari dua sisi. Kotak-kotak berubah
menjadi belah ketupat, dan hampir semua garis menjadi miring. Tetapi ada satu garis yang tetap
lurus dan tetap sejajar dengan dirinya sendiri: garis yang searah tarikan. Panjangnya berubah,
arahnya tidak. Kadang ada juga garis lain yang tidak miring tetapi terbalik arahnya, seperti
panah yang diputar setengah lingkaran.

Kalau kita hanya diberi tahu berapa kali setiap garis istimewa itu mulur, kita sudah tahu seluruh
tarikan, tanpa perlu menggambar jaringannya lagi. Itu cara yang sangat hemat: tiga atau empat
angka menggantikan ribuan titik pada lembaran.

Masalahnya, tidak semua tarikan begitu rapi. Coba pelintir sedikit lembarnya sambil menarik.
Sekarang tidak ada garis yang benar-benar lurus lagi, dan cara hemat tadi gagal total. Ada
tarikan yang seluruhnya bisa dijelaskan oleh beberapa angka mulur, ada yang tidak bisa.

Pertanyaannya: dari gambar tarikan saja, bisakah kita menebak apakah ada garis istimewa yang
tidak miring, dan berapa banyak? Di bab ini kita akan tahu mengapa sebagian tarikan punya garis
seperti itu dan sebagian tidak, mengapa ada tarikan yang seluruhnya dapat diganti oleh sekadar
beberapa angka, dan mengapa mesin hitung kadang salah menebak pada tarikan yang hampir pelintir.
\end{untukanak}

\section{Arah yang tidak berputar}
\tingkat{dasar}

Sebuah matriks persegi bekerja pada vektor: ia memutar, meregangkan, atau memipihkan. Hampir
semua vektor keluar dengan arah yang berbeda dari arah masuknya. Sebagian kecil tidak. Vektor
yang keluar searah dengan masuknya, hanya berbeda panjang, menandai arah yang dipertahankan
matriks itu. Arah semacam itu yang kita beri nama lebih dulu, sebelum ada alat hitung apa pun.

\kamus{nilai eigen}{eigenvalue}{Pengali pada arah yang tidak diputar: arahnya tetap, panjangnya
dikali angka itu.}

\begin{definisi}[Nilai eigen dan vektor eigen]
Misalkan $A$ matriks $n\times n$ dengan entri real. Skalar $\lambda$ disebut \emph{nilai eigen}
matriks $A$ bila ada vektor tak nol $\mathbf{v}\in\mathbb{R}^{n}$ dengan
$A\mathbf{v}=\lambda\mathbf{v}$. Vektor $\mathbf{v}$ disebut \emph{vektor eigen} untuk
$\lambda$.\index{nilai eigen}\index{vektor eigen}
\end{definisi}

Syarat $\mathbf{v}\neq\mathbf{0}$ bukan hiasan. Tanpa syarat itu, $\mathbf{0}$ memenuhi
$A\mathbf{0}=\lambda\mathbf{0}$ untuk setiap $\lambda$, sehingga istilahnya kehilangan arti. Yang
kita cari justru arah, dan titik nol tidak punya arah.

Persamaan $A\mathbf{v}=\lambda\mathbf{v}$ dapat ditulis ulang sebagai
$(A-\lambda I)\mathbf{v}=\mathbf{0}$. Jadi $\lambda$ adalah nilai eigen tepat ketika matriks
$A-\lambda I$ mempunyai inti yang tidak trivial, yaitu tepat ketika $A-\lambda I$ singular.
Karena itu pertanyaan tentang nilai eigen selalu dapat diterjemahkan menjadi pertanyaan tentang
kekurangan pangkat, dan Bab 8 sudah menyediakan alat untuk itu.

Sebelum menghitung, satu peringatan penting. Tidak setiap matriks real punya nilai eigen real.
Matriks putar
\[
  R=\begin{pmatrix}0&-1\\ 1&0\end{pmatrix}
\]
mengirim $(1,0)$ ke $(0,1)$. Setiap vektor real berputar sembilan puluh derajat, sehingga tidak ada
garis yang dipertahankan; $R^{2}=-I$ memutar setengah lingkaran setelah dua langkah.
Ketidakadaan nilai eigen real di sini bukan kegagalan hitung, melainkan pernyataan geometris.

Untuk mencari nilai eigen tanpa menebak, kita pakai determinan: nilai eigen adalah bilangan yang
membuat matriks kehilangan inversnya setelah diagonalnya dikurangi bilangan itu.

\kamus{polinomial karakteristik}{characteristic polynomial}{Daftar bilangan yang membuat matriks
menjadi singular, ditulis sebagai polinomial dalam satu peubah.}

\begin{definisi}[Polinomial karakteristik]
Untuk matriks $n\times n$, \emph{polinomial karakteristik} didefinisikan sebagai
$p_{A}(\lambda)=\det(\lambda I-A)$.\index{polinomial karakteristik}
\end{definisi}

Kita memakai $\lambda I-A$, bukan $A-\lambda I$. Keduanya berbeda faktor $(-1)^{n}$, jadi akarnya
sama, tetapi pilihan pertama membuat polinomialnya monik, dan itu memudahkan pembacaan koefisien
lain seperti terlihat pada akibat berikutnya.

\begin{teorema}
Skalar $\lambda$ adalah nilai eigen $A$ jika dan hanya jika $p_{A}(\lambda)=0$.
\end{teorema}

\begin{proof}
Misalkan $A\mathbf{v}=\lambda\mathbf{v}$ dengan $\mathbf{v}\neq\mathbf{0}$. Maka
$(A-\lambda I)\mathbf{v}=\mathbf{0}$, sehingga $A-\lambda I$ mempunyai inti tidak trivial dan
karena itu singular. Menurut Bab 8, singularitas setara dengan determinan nol, jadi
$\det(A-\lambda I)=0$, yang berbeda faktor $(-1)^{n}$ dari $\det(\lambda I-A)=0$.

Sebaliknya, bila $\det(\lambda I-A)=0$ maka $\lambda I-A$ singular. Ruang intinya tidak trivial,
jadi ada $\mathbf{v}\neq\mathbf{0}$ dengan $(\lambda I-A)\mathbf{v}=\mathbf{0}$, yaitu
$A\mathbf{v}=\lambda\mathbf{v}$.
\end{proof}

Polinomial $p_{A}$ berderajat $n$ dan monik, jadi ia punya paling banyak $n$ akar berbeda. Di
atas $\mathbb{C}$ ia selalu terurai menjadi $n$ faktor linear, sehingga setiap matriks persegi
kompleks punya tepat $n$ nilai eigen bila dihitung dengan multiplisitas. Di atas $\mathbb{R}$ hal
itu bisa gagal: $p_{R}(\lambda)=\lambda^{2}+1$ tidak punya akar real, tetapi punya sepasang akar
kompleks $\pm i$.

Dua koefisien teratas polinomial karakteristik memuat dua besaran yang sudah kita kenal.

\begin{akibat}
Dengan nilai eigen dihitung beserta multiplisitasnya,
$\operatorname{tr}(A)=\sum_{i}\lambda_{i}$ dan $\det(A)=\prod_{i}\lambda_{i}$.
\index{jejak}
\end{akibat}

\begin{proof}
Koefisien $\lambda^{n-1}$ pada $\det(\lambda I-A)$ adalah $-\operatorname{tr}(A)$. Suku itu hanya
dihasilkan oleh permutasi identitas pada rumus Leibniz, karena setiap permutasi lain kehilangan
paling sedikit dua entri diagonal. Suku konstannya adalah $p_{A}(0)=\det(-A)=(-1)^{n}\det(A)$.
Membandingkan kedua koefisien itu dengan bentuk
$\prod_{i}(\lambda-\lambda_{i})=\lambda^{n}-(\sum_{i}\lambda_{i})\lambda^{n-1}+\dots+(-1)^{n}
\prod_{i}\lambda_{i}$ memberi kedua identitas.
\end{proof}

Akibat itu memberi pemeriksaan murah: setelah menghitung nilai eigen dengan susah payah, jumlah dan
hasil kalinya harus cocok dengan jejak dan determinan. Bila tidak cocok, ada kesalahan aritmetika.

\begin{bacakeras}
Dua suara, bergantian.

Suara pertama: apa yang dilakukan sebuah matriks?\\
Suara kedua: memutar, meregangkan, memipihkan.\\
Suara pertama: kepada siapa?\\
Suara kedua: kepada hampir semua vektor.\\
Suara pertama: lalu ada yang tidak berputar?\\
Suara kedua: ada. Satu arah. Kadang dua. Tidak pernah lebih dari banyaknya kolom.\\
Suara pertama: panjangnya berubah?\\
Suara kedua: berubah. Arahnya tetap.\\
Suara pertama: berapa kali?\\
Suara kedua: itu angkanya. Angka itu nilai eigen.\\
Suara pertama: dan arahnya?\\
Suara kedua: itu vektor eigen. Satu arah, satu angka.\\
Suara pertama: kalau arahnya berbalik?\\
Suara kedua: angkanya negatif. Garisnya tetap sama, hanya menghadap ke belakang.\\
Suara pertama: kalau angkanya nol?\\
Suara kedua: arah itu hilang, dan matriks kehilangan inversnya.

Bersama: arah yang tidak berubah menyimpan angka yang mengubah panjangnya.

Ulangi satu kali dengan suara biasa: arah yang tidak berubah menyimpan angka yang mengubah panjangnya.
\end{bacakeras}

\begin{resep}
Untuk $A=\begin{pmatrix}a&b\\c&d\end{pmatrix}$ yang entrinya sudah diketahui:
\begin{enumerate}
  \item Kurangi $\lambda$ dari kedua entri diagonal, sehingga tersisa $A-\lambda I$.
  \item Hitung $\det(A-\lambda I)$, atau langsung pakai $\lambda^{2}-(a+d)\lambda+(ad-bc)$.
  \item Selesaikan persamaan kuadrat itu; periksa jumlah akarnya $a+d$ dan hasil kalinya $ad-bc$.
  \item Untuk setiap akar $\lambda$, selesaikan $(A-\lambda I)\mathbf{v}=\mathbf{0}$ dan ambil satu
        solusi tak nol.
\end{enumerate}
Langkah 3 dan 4 tidak boleh ditukar: vektor eigen baru punya arti setelah $\lambda$ diketahui, dan
urutan kolom $P$ harus mengikuti urutan entri diagonal $D$.
\end{resep}

Sekarang satu contoh dengan tangan. Ambil
\[
  A=\begin{pmatrix}4&2\\ 1&3\end{pmatrix},\qquad
  p_{A}(\lambda)=\lambda^{2}-7\lambda+10=(\lambda-5)(\lambda-2).
\]
Jejaknya $7=5+2$ dan determinannya $10=5\cdot2$, cocok. Untuk $\lambda=5$ kita selesaikan
$(A-5I)\mathbf{v}=\mathbf{0}$, yaitu $-v_{1}+2v_{2}=0$ dan $v_{1}-2v_{2}=0$; kedua baris
mengatakan hal yang sama, jadi $\mathbf{v}=(2,1)$. Untuk $\lambda=2$ diperoleh
$2v_{1}+2v_{2}=0$, jadi $\mathbf{v}=(1,-1)$. Pemeriksaan langsung:
\[
  A\begin{pmatrix}2\\ 1\end{pmatrix}=\begin{pmatrix}10\\ 5\end{pmatrix}=5
  \begin{pmatrix}2\\ 1\end{pmatrix},\qquad
  A\begin{pmatrix}1\\ -1\end{pmatrix}=\begin{pmatrix}2\\ -2\end{pmatrix}=2
  \begin{pmatrix}1\\ -1\end{pmatrix}.
\]
Kedua pemeriksaan itu dapat dilihat pada bidang. Gambar~\ref{fig:panah-eigen} menggambar kedua
arah eigen beserta bayangannya, dan satu vektor yang bukan vektor eigen sebagai pembanding.

\begin{figure}[htbp]
\centering
\begin{tikzpicture}[x=0.4cm,y=0.4cm,>=latex]
  \draw[->,layergrey!60] (-0.8,0) -- (11.4,0) node[right,font=\scriptsize] {$x$};
  \draw[->,layergrey!60] (0,-2.9) -- (0,6.4) node[above,font=\scriptsize] {$y$};
  \draw[->,thick,layergrey] (0,0) -- (1,0);
  \draw[->,thick,dashed,layergrey] (0,0) -- (4,1)
        node[right,font=\scriptsize] {$(4,1)$};
  \draw[->,very thick,layerteal] (0,0) -- (2,1)
        node[above left,font=\scriptsize] {$(2,1)$};
  \fill[layerteal] (2,1) circle (2pt);
  \draw[->,very thick,accent] (0,0) -- (10,5)
        node[above,font=\scriptsize] {$(10,5)=5\,(2,1)$};
  \draw[->,very thick,layerindigo] (0,0) -- (1,-1)
        node[below left,font=\scriptsize] {$(1,-1)$};
  \fill[layerindigo] (1,-1) circle (2pt);
  \draw[->,very thick,layerindigo] (0,0) -- (2,-2)
        node[below right,font=\scriptsize] {$(2,-2)=2\,(1,-1)$};
\end{tikzpicture}
\caption{Arah eigen matriks $A=\begin{pmatrix}4&2\\1&3\end{pmatrix}$. Vektor $(2,1)$ keluar
sebagai $(10,5)$, yaitu $5$ kali lipat tanpa berputar; vektor $(1,-1)$ keluar sebagai $(2,-2)$,
yaitu $2$ kali lipat. Vektor $(1,0)$ keluar sebagai $(4,1)$ dan berputar, jadi ia bukan vektor
eigen.}
\label{fig:panah-eigen}
\end{figure}

Vektor eigen hanya ditentukan sampai pengali: $(4,2)$ juga vektor eigen untuk $5$. Yang tunggal
adalah arahnya, bukan panjangnya. Kebebasan itu akan kita pakai untuk merapikan matriks nanti.

Untuk matriks $2\times2$ akar itu dapat ditulis langsung. Karena
$p_{A}(\lambda)=\lambda^{2}-\operatorname{tr}(A)\lambda+\det(A)$, rumus akar memberi
\[
  \lambda_{1,2}=\frac{\operatorname{tr}(A)\pm\sqrt{\operatorname{tr}(A)^{2}-4\det(A)}}{2},
\]
sehingga tanda $\Delta=\operatorname{tr}(A)^{2}-4\det(A)$ menentukan seluruh spektrumnya. Untuk $A$
di atas, $\Delta=49-40=9$, dan akarnya $5$ serta $2$. Bila $\Delta=0$ kedua akar berimpit, dan bila
$\Delta<0$ tidak ada akar real: matriks putar tadi memberi $\Delta=0-4=-4$.

\begin{figure}[htbp]
\centering
\begin{tikzpicture}[x=0.85cm,y=0.85cm,>=latex,font=\scriptsize]
  \draw[->,layergrey!60] (-1.2,0.35) -- (6.6,0.35);
  \foreach \x in {0,1,2,3,4,5,6} \draw[layergrey!50] (\x,0.26) -- (\x,0.44);
  \fill[layerteal] (2,0.35) circle (2.2pt);
  \fill[layerteal] (5,0.35) circle (2.2pt);
  \node[above=2pt] at (2,0.35) {$2$};
  \node[above=2pt] at (5,0.35) {$5$};
  \node[left,align=right,text width=3.0cm] at (-1.3,0.35) {$\Delta=9$: dua akar real};

  \draw[->,layergrey!60] (-1.2,-1.25) -- (6.6,-1.25);
  \foreach \x in {0,1,2,3,4,5,6} \draw[layergrey!50] (\x,-1.34) -- (\x,-1.16);
  \fill[layerorange] (1,-1.25) circle (2.2pt);
  \node[above=2pt] at (1,-1.2) {$1$};
  \node[left,align=right,text width=3.0cm] at (-1.3,-1.25) {$\Delta=0$: satu akar ganda};

  \draw[->,layergrey!60] (-1.2,-2.85) -- (6.6,-2.85);
  \foreach \x in {0,1,2,3,4,5,6} \draw[layergrey!50] (\x,-2.94) -- (\x,-2.76);
  \draw[layergrey!80] (0,-3.55) -- (0,-2.15);
  \draw[accent,fill=white,thick] (0,-2.3) circle (2.4pt);
  \draw[accent,fill=white,thick] (0,-3.4) circle (2.4pt);
  \node[left=2pt,accent] at (0,-2.3) {$i$};
  \node[left=2pt,accent] at (0,-3.4) {$-i$};
  \node[left,align=right,text width=3.0cm] at (-1.3,-2.85) {$\Delta=-4$: tidak ada akar real};
\end{tikzpicture}
\caption{Spektrum matriks $2\times2$ pada garis bilangan, dipisahkan oleh tanda
$\Delta=\operatorname{tr}(A)^{2}-4\det(A)$. Baris atas: $A$ dengan $\Delta=49-40=9$ memberi akar
$2$ dan $5$. Baris tengah: matriks guntingan dengan $\Delta=0$ memberi akar $1$ yang dihitung dua
kali. Baris bawah: matriks putar dengan $\Delta=0-4=-4$ tidak memotong garis real, dan kedua
akarnya, $\pm i$, berada di luar garis itu.}
\label{fig:spektrum}
\end{figure}

Gambar~\ref{fig:spektrum} menyusun ketiga kemungkinan itu pada satu garis bilangan: dua titik,
satu titik yang dihitung dua kali, atau tidak ada titik sama sekali. Yang terakhir bukan kekurangan
gambar, melainkan pernyataan bahwa garis real tidak memuat arah yang dipertahankan.

\begin{tangan}
Bahan: satu karet gelang, dua pensil, selembar kertas, penggaris.

Letakkan kedua pensil sejajar, lalu bentangkan karet gelang memutari keduanya sehingga terbentuk
persegi panjang. Di kertas, gambar tiga panah dari sudut kiri bawah persegi itu, masing-masing pada
sudut 30, 45, dan 60 derajat terhadap sisi bawah. Beri nama ketiganya.

Geser pensil kanan dua kali lebih jauh, sehingga karet meregang. Bandingkan setiap panah dengan
panjang aslinya. Panah yang sejajar sisi tarikan tetap pada arahnya; dua panah lain miring, dan
yang paling miring adalah yang paling jauh dari sisi tarikan.

Sekarang putar seluruh rangka sembilan puluh derajat dan ulangi. Arah yang bertahan ikut berputar
bersama karet, karena yang istimewa adalah arahnya terhadap karet, bukan terhadap meja.

Yang harus terlihat: satu arah bertahan, dua arah berputar, dan angka regangan pada arah yang
bertahan itulah pengali yang kita sebut nilai eigen.
\end{tangan}

\section{Ruang eigen dan dua multiplisitas}
\tingkat{menengah}

Satu nilai eigen biasanya disertai lebih dari satu vektor eigen. Himpunan seluruh vektor eigen
untuk satu nilai eigen, ditambah vektor nol, bukan kumpulan acak: ia subruang. Dimensinya nanti
menentukan apakah matriksnya dapat disederhanakan.

\begin{definisi}[Ruang eigen]
Untuk nilai eigen $\lambda$ dari $A$, \emph{ruang eigen} adalah
$E_{\lambda}=\operatorname{null}(A-\lambda I)=\{\mathbf{v}:A\mathbf{v}=\lambda\mathbf{v}\}$.
\index{ruang eigen}
\end{definisi}

\kamus{ruang eigen}{eigenspace}{Seluruh arah yang tidak diputar oleh matriks, dikumpulkan menjadi
satu subruang.}

Ruang eigen selalu subruang karena ia inti sebuah matriks (Bab 7). Dimensinya paling sedikit satu
tepat ketika $\lambda$ benar-benar nilai eigen.

Ruang eigen untuk nilai eigen berbeda hanya beririsan di titik nol.

\begin{proposisi}
Bila $\lambda\neq\mu$, maka $E_{\lambda}\cap E_{\mu}=\{\mathbf{0}\}$.
\end{proposisi}

\begin{proof}
Ambil $\mathbf{v}$ di irisan itu. Maka $A\mathbf{v}=\lambda\mathbf{v}$ dan
$A\mathbf{v}=\mu\mathbf{v}$, sehingga $(\lambda-\mu)\mathbf{v}=\mathbf{0}$. Karena
$\lambda-\mu\neq0$, harus $\mathbf{v}=\mathbf{0}$.
\end{proof}

Sifat itu naik ke sejumlah berhingga vektor eigen, dengan induksi. Hasilnya lemma yang menjadi
tulang punggung bab ini.

\begin{lemma}
Vektor eigen yang bersesuaian dengan nilai eigen berbeda saling bebas linear.
\end{lemma}

\begin{proof}
Kita buktikan dengan induksi pada banyaknya vektor. Untuk satu vektor pernyataannya benar karena
vektor eigen tidak nol. Andaikan pernyataan berlaku untuk $m-1$ vektor, dan misalkan
\[
  c_{1}\mathbf{v}_{1}+\dots+c_{m}\mathbf{v}_{m}=\mathbf{0}
\]
dengan $\lambda_{1},\dots,\lambda_{m}$ berbeda. Kalikan persamaan itu dengan $A$, dan pakai
$A\mathbf{v}_{i}=\lambda_{i}\mathbf{v}_{i}$:
\[
  c_{1}\lambda_{1}\mathbf{v}_{1}+\dots+c_{m}\lambda_{m}\mathbf{v}_{m}=\mathbf{0}.
\]
Kurangi $\lambda_{m}$ kali persamaan pertama dari persamaan kedua:
\[
  c_{1}(\lambda_{1}-\lambda_{m})\mathbf{v}_{1}+\dots
  +c_{m-1}(\lambda_{m-1}-\lambda_{m})\mathbf{v}_{m-1}=\mathbf{0}.
\]
Menurut hipotesis induksi, $\mathbf{v}_{1},\dots,\mathbf{v}_{m-1}$ bebas, jadi
$c_{i}(\lambda_{i}-\lambda_{m})=0$ untuk $i<m$. Karena $\lambda_{i}\neq\lambda_{m}$, semua
$c_{i}=0$ untuk $i<m$. Persamaan semula tinggal $c_{m}\mathbf{v}_{m}=\mathbf{0}$, sehingga
$c_{m}=0$.
\end{proof}

Lemma itu memberi batas atas yang rapi: matriks $n\times n$ tidak mungkin punya lebih dari $n$
vektor eigen bebas. Kapan batas itu tercapai?

Untuk menjawabnya kita butuh dua cara menghitung banyaknya kemunculan satu nilai eigen: dari
polinomial dan dari dimensi ruang eigen. Keduanya sering berbeda, dan selisihnya yang menentukan.

\kamus{multiplisitas}{multiplicity}{Berapa kali sebuah nilai eigen muncul, dihitung dengan dua cara
yang tidak selalu sama.}

\begin{definisi}[Multiplisitas aljabar dan geometrik]
Misalkan $\lambda$ nilai eigen $A$. \emph{Multiplisitas aljabar} $\lambda$ adalah pangkat faktor
$(\lambda'-\lambda)$ pada $p_{A}(\lambda')$. \emph{Multiplisitas geometrik} $\lambda$ adalah
$\dim E_{\lambda}=\operatorname{nullity}(A-\lambda I)$.\index{multiplisitas}
\end{definisi}

Multiplisitas geometrik mudah dihitung: kurangi $\lambda$ dari diagonal, lalu hitung nullity
matriks hasilnya. Multiplisitas aljabar menuntut pemfaktoran polinomial. Hubungan keduanya selalu
satu arah.

\begin{lemma}
Untuk setiap nilai eigen $\lambda$ berlaku
$1\le\dim E_{\lambda}\le$ multiplisitas aljabar $\lambda$.
\end{lemma}

\begin{proof}
Pertidaksamaan kiri benar karena $\lambda$ nilai eigen, jadi $E_{\lambda}$ memuat vektor tak nol.

Untuk pertidaksamaan kanan, tulis $k=\dim E_{\lambda}$ dan pilih basis
$\mathbf{u}_{1},\dots,\mathbf{u}_{k}$ dari $E_{\lambda}$. Lengkapi menjadi basis
$\mathbf{u}_{1},\dots,\mathbf{u}_{n}$ dari $\mathbb{R}^{n}$, yang mungkin menurut Bab 7. Dalam
basis itu, kolom ke-$j$ matriks $A$ untuk $j\le k$ adalah $A\mathbf{u}_{j}=\lambda\mathbf{u}_{j}$,
yaitu $\lambda$ pada posisi diagonal dan nol di tempat lain. Jadi matriks $A$ dalam basis tersebut
berbentuk
\[
  \begin{pmatrix}\lambda I_{k}&C\\ 0&D\end{pmatrix}
\]
dengan $C$ berukuran $k\times(n-k)$ dan $D$ berukuran $(n-k)\times(n-k)$. Determinan
$\lambda' I-A$ dalam basis itu, lewat ekspansi kofaktor pada $k$ kolom pertama atau lewat rumus
determinan matriks blok segitiga, sama dengan $(\lambda'-\lambda)^{k}\det(\lambda' I-D)$. Jadi
$(\lambda'-\lambda)^{k}$ membagi $p_{A}$, sehingga multiplisitas aljabar $\lambda$ paling sedikit
$k$.
\end{proof}

Sekarang lihat kedua kemungkinan pada satu matriks. Ambil
\[
  A=\begin{pmatrix}2&1&1\\ 1&2&1\\ 1&1&2\end{pmatrix},\qquad
  p_{A}(\lambda)=(\lambda-1)^{2}(\lambda-4).
\]
Nilai eigennya $1$ dan $4$, dengan multiplisitas aljabar $2$ dan $1$. Jejaknya $6=4+1+1$ dan
determinannya $4=4\cdot1\cdot1$, cocok. Untuk $\lambda=4$, persamaan $(A-4I)\mathbf{v}=
\mathbf{0}$ memberi $v_{1}=v_{2}=v_{3}$, jadi $E_{4}$ direntang $(1,1,1)$. Untuk $\lambda=1$,
matriks $A-I$ adalah matriks yang seluruh entrinya $1$, sehingga intinya himpunan vektor yang
jumlah entrinya nol. Himpunan itu berdimensi dua, dengan basis $(1,-1,0)$ dan $(1,0,-1)$. Jadi
multiplisitas geometrik $1$ adalah $2$, sama dengan multiplisitas aljabarnya.

Nilai eigen berulang tidak otomatis menjadi masalah; yang menjadi masalah ruang eigen yang terlalu
kecil, dan itu terjadi pada matriks berikut.

\section{Kapan matriks dapat didiagonalkan}
\tingkat{menengah}

Kalau boleh memilih basis, kita ingin matriksnya sesederhana mungkin: diagonal, sehingga setiap
vektor basis hanya diskalakan dan tidak dicampur dengan yang lain. Apakah selalu ada basis semacam
itu?

\begin{definisi}[Dapat didiagonalkan]
Matriks $n\times n$ disebut \emph{dapat didiagonalkan} bila ada matriks $P$ yang invertible dan
matriks diagonal $D$ dengan $P^{-1}AP=D$, atau setara dengan itu $A=PDP^{-1}$.
\end{definisi}

\kamus{diagonalisasi}{diagonalization}{Mengganti alat ukur sampai seluruh pekerjaan matriks
tinggal beberapa perkalian angka.}

Kolom-kolom $P$ nanti akan berisi vektor eigen, dan $D$ berisi nilai eigen pada diagonalnya.
Karena itu pertanyaan tentang diagonalisasi sebenarnya pertanyaan tentang ketersediaan vektor
eigen yang cukup banyak.

\begin{teorema}
Matriks $n\times n$ dapat didiagonalkan jika dan hanya jika ia mempunyai $n$ vektor eigen yang
bebas linear.\index{diagonalisasi}
\end{teorema}

\begin{proof}
Misalkan $P^{-1}AP=D$ dengan $D=\operatorname{diag}(d_{1},\dots,d_{n})$. Kalikan dari kiri
dengan $P$: $AP=PD$. Kolom ke-$j$ kedua ruas memberi
$A\mathbf{p}_{j}=d_{j}\mathbf{p}_{j}$, dengan $\mathbf{p}_{j}$ kolom ke-$j$ dari $P$. Karena $P$
invertible, setiap kolomnya tak nol, jadi $\mathbf{p}_{j}$ vektor eigen. Kolom-kolom $P$ juga
bebas linear karena $P$ invertible. Jadi $A$ punya $n$ vektor eigen bebas.

Sebaliknya, misalkan $\mathbf{p}_{1},\dots,\mathbf{p}_{n}$ vektor eigen bebas dengan
$A\mathbf{p}_{j}=\lambda_{j}\mathbf{p}_{j}$. Susun $P$ dari vektor-vektor itu sebagai kolom dan
$D=\operatorname{diag}(\lambda_{1},\dots,\lambda_{n})$. Maka $AP=PD$, dan karena $P$ invertible
kita peroleh $P^{-1}AP=D$.
\end{proof}

Akibat yang paling sering dipakai dalam hitungan tangan: banyaknya nilai eigen berbeda sudah cukup.

\begin{akibat}
Bila matriks $n\times n$ mempunyai $n$ nilai eigen berbeda, ia dapat didiagonalkan.
\end{akibat}

\begin{proof}
Ambil satu vektor eigen untuk setiap nilai eigen. Menurut lemma sebelumnya, $n$ vektor eigen
dengan $n$ nilai eigen berbeda saling bebas. Teorema di atas lalu memberi diagonalisasi.
\end{proof}

Kebalikannya tidak berlaku: contoh $3\times3$ tadi hanya punya dua nilai eigen berbeda, tetapi
tetap dapat didiagonalkan karena ruang eigen untuk $\lambda=1$ berdimensi dua. Syarat yang tepat
bukan banyaknya nilai eigen, melainkan jumlah dimensi ruang eigennya.

\begin{teorema}[Syarat multiplisitas]
Matriks $n\times n$ dapat didiagonalkan jika dan hanya jika untuk setiap nilai eigennya,
multiplisitas aljabar sama dengan multiplisitas geometrik. Syarat itu setara dengan
$\sum_{\lambda}\dim E_{\lambda}=n$.
\end{teorema}

\begin{proof}
Pertama, anggap jumlah dimensi seluruh ruang eigen sama dengan $n$. Untuk setiap $\lambda$ pilih
basis $E_{\lambda}$, lalu gabungkan seluruh basis itu menjadi daftar $\mathbf{v}_{1},\dots,
\mathbf{v}_{n}$. Daftar itu bebas linear: bila suatu kombinasi linear bernilai nol, kelompokkan
suku menurut nilai eigennya menjadi $\mathbf{u}_{1}+\dots+\mathbf{u}_{m}=\mathbf{0}$ dengan
$\mathbf{u}_{k}\in E_{\lambda_{k}}$. Menurut lemma tentang vektor eigen berbeda, setiap
$\mathbf{u}_{k}=\mathbf{0}$; karena setiap $\mathbf{u}_{k}$ kombinasi basis $E_{\lambda_{k}}$,
semua koefisiennya nol. Daftar itu juga berisi $n$ vektor di ruang berdimensi $n$, jadi ia basis
vektor eigen, dan teorema sebelumnya menutup jalan.

Kedua, anggap $A$ dapat didiagonalkan, jadi ada basis vektor eigen
$\mathbf{p}_{1},\dots,\mathbf{p}_{n}$. Dalam basis itu matriks $A$ diagonal, sehingga
$p_{A}(\lambda')=\prod_{j}(\lambda'-\lambda_{j})$ dan multiplisitas aljabar setiap nilai eigen
sama dengan banyaknya $j$ dengan $\lambda_{j}=\lambda$. Banyaknya $j$ itu tidak melebihi
$\dim E_{\lambda}$, karena vektor-vektor tersebut berada di dalam $E_{\lambda}$ dan bebas. Jadi
multiplisitas aljabar tidak melebihi geometrik. Lemma sebelumnya memberi arah sebaliknya, sehingga
keduanya sama. Terakhir, $\sum_{\lambda}\dim E_{\lambda}=n$ karena basis vektor eigen mengisi
seluruh ruang.
\end{proof}

Sekarang contoh yang gagal. Ambil
\[
  S=\begin{pmatrix}1&1\\ 0&1\end{pmatrix},\qquad p_{S}(\lambda)=(\lambda-1)^{2}.
\]
Satu-satunya nilai eigen adalah $1$, dengan multiplisitas aljabar $2$. Tetapi
$S-I=\begin{pmatrix}0&1\\ 0&0\end{pmatrix}$ berpangkat satu, jadi
$\dim E_{1}=2-1=1$. Multiplisitas geometriknya kurang satu, dan menurut teorema di atas $S$ tidak
dapat didiagonalkan.

Secara geometris $S$ adalah guntingan: ia mempertahankan arah $(1,0)$ dan menggeser $(0,1)$ ke
$(1,1)$. Tidak ada basis yang membuat peta ini sekadar penskalaan, karena penskalaan tidak pernah
menggeser sesuatu yang tadinya sejajar. Kegagalan itu sifat petanya, bukan pilihan basis.

\begin{untukanak}
Pernahkah kamu mencoba merapikan sesuatu yang ternyata tidak bisa dirapikan? Ambil satu tumpukan
kartu, lalu geser bagian atasnya sedikit ke samping. Gambar pada setiap kartu tetap sama, tetapi
setiap titik bergeser sejauh jaraknya dari alas. Tidak ada cara memutar tumpukan itu supaya
geseran tersebut berubah menjadi sekadar memanjangkan atau memendekkan sesuatu. Geseran tetap
geseran.

Ada benda lain yang lebih ramah. Karet yang ditarik lurus dapat dijelaskan hanya dengan berapa
kali setiap arah mulur. Kalau kita tahu angka itu untuk beberapa arah penting, kita tahu seluruh
tarikan, termasuk untuk titik yang tidak pernah kita ukur.

Pertanyaannya: bagaimana membedakan keduanya hanya dari gambarnya, tanpa mencoba semua
kemungkinan? Di bab ini kita akan tahu mengapa geseran selalu menyimpan satu arah yang tidak bisa
dirapikan, mengapa ada tarikan yang punya arah istimewa cukup banyak, dan mengapa beda itu
menentukan apakah pangkat kesepuluh sebuah mesin hitung masih dapat dipercaya.
\end{untukanak}

Guntingan mudah dikenali karena segitiga. Kegagalan yang sama muncul pada matriks yang tidak tampak
segitiga. Ambil
\[
  B=\begin{pmatrix}-1&4&0\\ -1&3&0\\ 1&-2&2\end{pmatrix},\qquad
  p_{B}(\lambda)=(\lambda-1)^{2}(\lambda-2).
\]
Untuk $\lambda=1$, matriks $B-I$ berpangkat dua, jadi $\dim E_{1}=3-2=1$, padahal
multiplisitas aljabarnya $2$. Satu-satunya vektor eigen untuk $1$ adalah $(2,1,0)$. Untuk
$\lambda=2$ ruang eigennya direntang $(0,0,1)$ dan berdimensi satu, sesuai multiplisitas
aljabarnya. Jumlah dimensi seluruh ruang eigen hanya $2<3$, jadi $B$ tidak dapat didiagonalkan.
Matriks semacam ini disebut \emph{defective}: ia kekurangan vektor eigen karena ruang eigennya
sempit.\index{defective}

\kamus{defective}{defective}{Matriks yang nilai eigennya berulang tetapi ruang eigennya terlalu
sempit untuk menjadi basis.}

Satu catatan sebelum kita lanjut. Dapat atau tidaknya didiagonalkan bergantung pada medan
bilangannya: $R$ tidak dapat didiagonalkan atas $\mathbb{R}$, tetapi atas $\mathbb{C}$ ia punya dua
nilai eigen berbeda, $\pm i$, sehingga dapat didiagonalkan. Kalimat tanpa keterangan medan karena
itu tidak lengkap. Di sisa bab ini medan yang dipakai adalah $\mathbb{R}$ kecuali disebut lain.

\section{Matriks serupa dan invarian}
\tingkat{lanjut}

Diagonalisasi adalah perubahan basis yang menghasilkan bentuk paling sederhana. Tetapi tidak
setiap matriks dapat mencapai bentuk itu, dan pertanyaan berikutnya wajar: apa yang tetap sama
pada semua matriks yang dihubungkan oleh perubahan basis?\index{invarian}

\begin{definisi}[Matriks serupa]
Dua matriks $n\times n$, $A$ dan $B$, disebut \emph{serupa} bila ada matriks invertible $P$
dengan $B=P^{-1}AP$.\index{serupa}
\end{definisi}

\kamus{serupa}{similar}{Dua laporan tentang peta yang sama dari dua ruang ukur yang berbeda.}

Keserupaan adalah relasi ekuivalen: setiap matriks serupa dengan dirinya sendiri, relasinya
simetris karena $(P^{-1})^{-1}BP^{-1}=A$, dan transitif karena
$(PQ)^{-1}A(PQ)=Q^{-1}(P^{-1}AP)Q$. Kelas-kelas itu terbagi rapi, dan diagonalisasi berarti
menanyakan apakah suatu kelas memuat matriks diagonal.

\begin{proposisi}
Bila $A$ dan $B$ serupa, maka $p_{A}=p_{B}$. Akibatnya $A$ dan $B$ mempunyai nilai eigen yang
sama beserta multiplisitas aljabarnya, dan juga jejak, determinan, pangkat, serta nullity yang
sama.
\end{proposisi}

\begin{proof}
Misalkan $B=P^{-1}AP$. Maka
\[
  \lambda I-B=\lambda I-P^{-1}AP=P^{-1}(\lambda I-A)P,
\]
sehingga
$\det(\lambda I-B)=\det(P^{-1})\det(\lambda I-A)\det(P)=\det(\lambda I-A)$ karena
$\det(P^{-1})\det(P)=1$. Jadi polinomial karakteristiknya sama, dan semua data yang diturunkan
dari koefisiennya ikut sama. Kesamaan pangkat dan nullity diperoleh dari $B-\lambda I=
P^{-1}(A-\lambda I)P$: perkalian dengan matriks invertible tidak mengubah pangkat, dan peringkat
itulah yang menentukan nullity lewat rank-nullity.
\end{proof}

Kebalikannya tidak berlaku, dan contoh penyangkalnya sederhana: $I_{2}$ dan guntingan $S$ punya
polinomial karakteristik sama, $(\lambda-1)^{2}$, tetapi tidak serupa. Bila $P^{-1}I_{2}P=S$, ruas
kiri sama dengan $I_{2}$ untuk setiap $P$. Jadi polinomial karakteristik bukan invarian lengkap; ia
hanya menyaring.

Yang berubah dan yang tidak berubah dapat dipisahkan dengan rapi. Nilai eigen tidak berubah.
Vektor eigennya berubah, dan berubah dengan cara yang mudah diingat.

\begin{proposisi}
Bila $B=P^{-1}AP$ dan $A\mathbf{v}=\lambda\mathbf{v}$, maka
$B(P^{-1}\mathbf{v})=\lambda(P^{-1}\mathbf{v})$.
\end{proposisi}

\begin{proof}
Hitung $B(P^{-1}\mathbf{v})=P^{-1}APP^{-1}\mathbf{v}=P^{-1}A\mathbf{v}
=P^{-1}(\lambda\mathbf{v})=\lambda(P^{-1}\mathbf{v})$.
\end{proof}

Jadi vektor eigen ikut berpindah bersama koordinatnya, sedangkan pengalinya tinggal. Itu masuk
akal: arah yang dipertahankan adalah sifat peta, sedangkan daftar koordinatnya bergantung pada
basis. Karena itu nilai eigen dipakai sebagai sidik jari matriks, sedangkan vektor eigen tidak.

Di sisi lain, kebebasan memilih $P$ menjelaskan mengapa dua diagonalisasi dapat tampak berbeda
padahal keduanya benar: kolom $P$ boleh dikalikan skalar bukan nol atau dipermutasi asalkan $D$ ikut
berpermutasi. Yang harus sama adalah himpunan nilai eigen beserta multiplisitasnya.

\section{Contoh hitung}
\tingkat{menengah}

Semua angka di bagian ini dihitung dengan tangan lebih dulu, lalu diperiksa dengan Python pada
bagian lab. Kita mulai dari matriks $3\times3$ yang sudah muncul pada pembahasan ruang eigen,
\[
  A=\begin{pmatrix}2&1&1\\ 1&2&1\\ 1&1&2\end{pmatrix}.
\]
Polinomial karakteristiknya dihitung dari $\lambda I-A$, yang seluruh entri diagonalnya
$\lambda-2$ dan seluruh entri lainnya $-1$. Dengan menulis $t=\lambda-2$ dan memakai ekspansi
sepanjang baris pertama,
\[
  p_{A}(\lambda)=t\,(t^{2}-1)-(-1)\,(-t-1)+(-1)\,(1+t)=t^{3}-3t-2 .
\]
Bentuk itu terurai karena $t=-1$ akarnya: $t^{3}-3t-2=(t+1)^{2}(t-2)$. Kembali ke $\lambda$,
\[
  p_{A}(\lambda)=(\lambda-1)^{2}(\lambda-4)=\lambda^{3}-6\lambda^{2}+9\lambda-4 .
\]
Jejaknya $6$ dan determinannya $4$, cocok dengan $4+1+1$ dan $4\cdot1\cdot1$.

Vektor eigen untuk $\lambda=4$: matriks $A-4I$ mempunyai entri diagonal $-2$ dan entri lain $1$.
Baris-barisnya mengatakan $-2v_{1}+v_{2}+v_{3}=0$ dan seterusnya; ketiga persamaan itu terpenuhi
tepat ketika $v_{1}=v_{2}=v_{3}$. Jadi $E_{4}$ direntang $(1,1,1)$.

Vektor eigen untuk $\lambda=1$: matriks $A-I$ seluruhnya berisi $1$, jadi intinya adalah himpunan
vektor dengan $v_{1}+v_{2}+v_{3}=0$. Basisnya dua vektor, misalnya $(1,-1,0)$ dan $(1,0,-1)$.
Dengan demikian
\[
  P=\begin{pmatrix}1&1&1\\ 1&-1&0\\ 1&0&-1\end{pmatrix},\qquad
  P^{-1}=\frac{1}{3}\begin{pmatrix}1&1&1\\ 1&-2&1\\ 1&1&-2\end{pmatrix},\qquad
  D=\operatorname{diag}(4,1,1).
\]
Determinan $P$ adalah $3$, dan itulah penyebut pada $P^{-1}$. Satu entri kita periksa dengan
tangan: entri baris pertama kolom pertama $PDP^{-1}$ adalah
$\tfrac{1}{3}(4\cdot1\cdot1+1\cdot1\cdot1+1\cdot1\cdot1)=2$, sesuai entri $A$. Pemeriksaan
menyeluruh dengan SymPy memberi $PDP^{-1}-A=0$ persis.

Sekarang pangkat kesepuluh.\index{pangkat matriks} Mengalikan sembilan kali tidak memberi gambaran
apa pun; diagonalisasi memberinya. Karena $A=I+J$ dengan $J$ matriks yang seluruh entrinya $1$, dan
$J^{2}=3J$, setiap pangkat $J$ kembali menjadi kelipatan $J$:
\[
  J^{2}=3J,\quad J^{3}=9J,\quad\dots,\quad J^{k}=3^{k-1}J .
\]
Maka
\[
  A^{k}=(I+J)^{k}=I+\Bigl(\sum_{j=1}^{k}\binom{k}{j}3^{j-1}\Bigr)J
  =I+\frac{4^{k}-1}{3}\,J,
\]
dengan langkah terakhir memakai rumus deret geometri
$\sum_{j=1}^{k}\binom{k}{j}3^{j-1}=\bigl((1+3)^{k}-1\bigr)/3$. Untuk $k=10$ kita punya
$4^{10}=1048576$ dan $(1048576-1)/3=349525$, sehingga
\[
  A^{10}=I+349525\,J=
  \begin{pmatrix}349526&349525&349525\\ 349525&349526&349525\\ 349525&349525&349526\end{pmatrix}.
\]
Setiap entrinya dapat diperiksa: diagonal $1+349525$, luar diagonal $349525$. Perkalian langsung
dengan Python memberi matriks yang sama, dengan selisih terbesar $1.05\times10^{-9}$ karena
aritmetika floating-point.

Bentuk rumus itu menjelaskan apa yang tidak terlihat pada hasil akhir. Suku $4^{k}$ menguasai arah
$(1,1,1)$; suku konstan $1$ bertahan pada arah yang jumlah entrinya nol. Nilai eigen $4$ dan $1$
adalah dua laju pertumbuhan, dan pangkat matriks hanyalah campuran keduanya. Bila matriks tidak
dapat didiagonalkan, campuran itu berubah bentuk: muncul $k\lambda^{k}$ atau, untuk $\lambda=1$,
faktor $k$ murni.

Contoh defektif tadi memperlihatkan hal itu. Untuk $B$ yang gagal didiagonalkan tadi,
\[
  B^{k}=
  \begin{pmatrix}1-2k&4k&0\\ -k&1+2k&0\\ 2^{k}-1&-2(2^{k}-1)&2^{k}\end{pmatrix},
\]
yang dapat dibuktikan dengan induksi dan cocok untuk $k=1$. Pada $k=10$ rumus itu memberi
\[
  B^{10}=\begin{pmatrix}-19&40&0\\ -10&21&0\\ 1023&-2046&1024\end{pmatrix},
\]
sama dengan hasil perkalian berulang. Entri kiri atas, $-19=1-2\cdot10$, tumbuh linear, bukan
eksponensial: nilai eigen $1$ mengendalikan arah itu, dan kekurangan vektor eigen menambahkan satu
faktor $k$.

Sekarang penerapan pertama: proses yang menetap. Untuk matriks rantai Markov $M$ di pembuka,
\index{rantai Markov}
\[
  p_{M}(\lambda)=\lambda^{3}-\tfrac{19}{10}\lambda^{2}+\tfrac{11}{10}\lambda-\tfrac15
  =(\lambda-1)\bigl(\lambda-\tfrac12\bigr)\bigl(\lambda-\tfrac25\bigr).
\]
Jejaknya $19/10$ dan determinannya $1/5$, cocok dengan $1+\tfrac12+\tfrac25$ dan
$1\cdot\tfrac12\cdot\tfrac25$. Vektor eigen untuk $\lambda=1$ dicari dari $M\mathbf{v}=\mathbf{v}$,
yaitu $(M-I)\mathbf{v}=\mathbf{0}$, ditambah syarat jumlah entrinya satu. Penyelesaiannya memberi
$\boldsymbol{\pi}=(17/30,\;4/15,\;1/6)$. Pemeriksaan satu koordinat dengan tangan:
\[
  \tfrac{8}{10}\cdot\tfrac{17}{30}+\tfrac{3}{10}\cdot\tfrac{8}{30}
  +\tfrac{2}{10}\cdot\tfrac{5}{30}
  =\tfrac{13.6+2.4+1.0}{30}=\tfrac{17}{30}.
\]
Jadi $\boldsymbol{\pi}$ benar-benar titik tetap.

Karena $M$ punya tiga nilai eigen berbeda, ia dapat didiagonalkan. Setiap titik awal karena itu
terurai sebagai
\[
  \mathbf{x}_{k}=\boldsymbol{\pi}+c_{2}\bigl(\tfrac12\bigr)^{k}\mathbf{v}_{2}
  +c_{3}\bigl(\tfrac25\bigr)^{k}\mathbf{v}_{3},
\]
dengan $c_{2},c_{3}$ bergantung pada $\mathbf{x}_{0}$. Suku kedua dan ketiga mengecil, dan yang
menentukan laju adalah $1/2$. Untuk $\mathbf{x}_{0}=(1,0,0)$ kita peroleh $c_{2}\mathbf{v}_{2}$
berorde $10^{-1}$, sehingga galat setelah sepuluh langkah berorde
$10^{-1}\cdot(1/2)^{10}\approx10^{-4}$, dan setelah empat puluh langkah berorde
$(1/2)^{40}\approx9.1\times10^{-13}$. Hasil hitung sesuai: $\mathbf{x}_{10}$ berbeda dari
$\boldsymbol{\pi}$ pada angka keempat, sedangkan $\mathbf{x}_{40}$ sudah sama sampai enam angka.

Penerapan kedua: sistem diferensial. Bila $\mathbf{x}(t)$ memenuhi
$\mathbf{x}'=A\mathbf{x}$ dengan $A$ di atas, solusinya $\mathbf{x}(t)=e^{At}\mathbf{x}(0)$.
Karena $A=I+J$ dan $J^{2}=3J$,
\[
  e^{At}=e^{t}e^{Jt}=e^{t}\Bigl(I+\frac{e^{3t}-1}{3}J\Bigr)
  =e^{t}I+\frac{e^{4t}-e^{t}}{3}J .
\]
Pada $t=0$ rumus itu memberi $I$. Untuk $t=0.5$ entri luar diagonalnya $(e^{2}-e^{0.5})/3=1.9134$,
dan Python memberi angka yang sama sampai empat angka. Komponen sepanjang $(1,1,1)$ tumbuh seperti
$e^{4t}$; komponen pada bidang jumlah nol tumbuh seperti $e^{t}$, sehingga solusi cepat sejajar
dengan $(1,1,1)$. Nilai eigen real positif berarti pertumbuhan; nilai eigen imajiner murni seperti
$\pm i$ berarti osilasi tanpa pertumbuhan.

Penerapan ketiga: iterasi pangkat. Mengalikan berulang lalu menormalkan hasilnya,
$\mathbf{x}_{k+1}=A\mathbf{x}_{k}/\lVert A\mathbf{x}_{k}\rVert$, membuat arah berputar menuju vektor
eigen dominan, karena $\lambda_{1}^{k}$ mengalahkan $\lambda_{2}^{k}$ bila
$|\lambda_{1}|>|\lambda_{2}|$. Untuk $A$ di atas mulai dari $(1,0)$, sudutnya berjalan $0^\circ$,
$14.04^\circ$, $21.25^\circ$, $24.39^\circ$, $25.69^\circ$, $26.21^\circ$, dan menuju
$\arctan(1/2)=26.57^\circ$; nisbah Rayleigh $\mathbf{x}^{\top}A\mathbf{x}/\mathbf{x}^{\top}
\mathbf{x}$ berjalan $4.00$, $4.65$, $4.88$, $4.96$, $4.98$, $4.99$, menuju $5$.
Gambar~\ref{fig:iterasi-pangkat} memperlihatkan putaran itu.

\begin{figure}[htbp]
\centering
\begin{tikzpicture}[scale=3.2,>=latex]
  \draw[layergrey!35] (0,0) circle (1);
  \draw[->,layergrey!60] (-1.25,0) -- (1.55,0) node[right,font=\scriptsize] {$x$};
  \draw[->,layergrey!60] (0,-0.35) -- (0,1.2) node[above,font=\scriptsize] {$y$};
  \draw[->,semithick,layergrey] (0,0) -- (0:1);
  \draw[->,thick,layerindigo!45] (0,0) -- (14.04:1);
  \draw[->,thick,layerindigo!70] (0,0) -- (21.25:1);
  \draw[->,very thick,layerindigo] (0,0) -- (24.39:1);
  \foreach \a in {0,14.04,21.25,24.39} \fill[layerindigo!70] (\a:1) circle (1.3pt);
  \draw[->,very thick,dashed,accent] (0,0) -- (26.565:1.42);
  \node[right,font=\scriptsize,accent,align=left] at (26.565:1.42)
        {$(2,1)$\\[1pt]$26.57^\circ$};
\end{tikzpicture}
\caption{Iterasi pangkat pada $A=\begin{pmatrix}4&2\\1&3\end{pmatrix}$ mulai dari $(1,0)$.
Panah penuh adalah arah setelah $k=0,1,2,3$ langkah, dengan sudut $0^\circ$, $14.04^\circ$,
$21.25^\circ$, dan $24.39^\circ$. Panah putus-putus adalah arah eigen dominan $(2,1)$ pada
$26.57^\circ$, yang didekati tanpa pernah dilewati; dua panah terakhir berimpit karena sudutnya
makin rapat, dan itu justru inti gambarnya.}
\label{fig:iterasi-pangkat}
\end{figure}

\kertas{Untuk $J$ matriks yang seluruh entrinya $1$, ambil $B=J-I$. Tentukan nilai eigen $B$
lebih dulu, lalu ramalkan $B^{10}$ memakai pola $I+\frac{\mu^{10}-(-1)^{10}}{3}J$ dengan $\mu$
nilai eigen bukan nol dari $B$. Periksa ramalan itu dengan Python, lalu jelaskan mengapa $B$ dapat
didiagonalkan meskipun punya nilai eigen berulang. Sebagai latihan tambahan, hitung $e^{Bt}$ dan
tentukan limit $t\to\infty$ dari $e^{-2t}e^{Bt}$.}

\section{Lab komputasi}
\tingkat{lanjut}

\begin{lab}
\begin{verbatim}
import numpy as np
from scipy.linalg import expm

np.set_printoptions(precision=4, suppress=True)

A = np.array([[2., 1., 1.],
              [1., 2., 1.],
              [1., 1., 2.]])
lam, V = np.linalg.eig(A)
print("nilai eigen A :", np.round(np.sort(lam.real), 10))
for i in range(3):
    sisa = A @ V[:, i] - lam[i] * V[:, i]
    print("  sisa Av-lv  :", round(float(np.linalg.norm(sisa)), 16))
print("A^10 lewat diagonalisasi :")
print(np.round((V @ np.diag(lam ** 10) @ np.linalg.inv(V)).real).astype(int))
print("A^10 lewat pangkat biasa :")
print(np.linalg.matrix_power(A, 10).astype(int))
print("selisih maks :", np.abs((V @ np.diag(lam ** 10)
                              @ np.linalg.inv(V)).real
                             - np.linalg.matrix_power(A, 10)).max())
print("cond V       :", round(float(np.linalg.cond(V)), 4))

A2 = np.array([[4., 2.],
               [1., 3.]])
x = np.array([1., 0.])
print("iterasi pangkat pada A2:")
for k in range(6):
    sudut = np.degrees(np.arctan2(x[1], x[0]))
    rq = float(x @ A2 @ x / (x @ x))
    print("  k=%d sudut=%7.2f rayleigh=%.4f" % (k, sudut, rq))
    x = A2 @ x / np.linalg.norm(A2 @ x)
print("sudut batas  :", round(float(np.degrees(np.arctan2(1., 2.))), 2))

B = np.array([[-1., 4., 0.],
              [-1., 3., 0.],
              [1., -2., 2.]])
print("eigen B      :", np.round(np.sort(np.linalg.eigvals(B).real), 10))
print("rank(B - I)  :", np.linalg.matrix_rank(B - np.eye(3)))
print("B^10         :")
print(np.linalg.matrix_power(B, 10).astype(int))

M = np.array([[0.8, 0.3, 0.2],
              [0.1, 0.6, 0.3],
              [0.1, 0.1, 0.5]])
print("jumlah kolom :", M.sum(axis=0))
print("eigen M      :", np.round(np.sort(np.linalg.eigvals(M).real), 10))
x = np.array([1., 0., 0.])
for _ in range(10):
    x = M @ x
print("x setelah 10 :", np.round(x, 6))
for _ in range(30):
    x = M @ x
print("x setelah 40 :", np.round(x, 6))

C = np.array([[1., 1000.], [0., 2.]])
Cp = C.copy()
Cp[1, 0] = 1e-6
print("eigen C      :", np.round(np.sort(np.linalg.eigvals(C).real), 6))
print("eigen C+1e-6 :", np.round(np.sort(np.linalg.eigvals(Cp).real), 6))
print("cond V(C)    :", round(float(np.linalg.cond(np.linalg.eig(C)[1])), 1))
print("expm(0.5A)   :", round(float(expm(0.5 * A)[0, 1]), 4))
print("rumus        :", round(float((np.exp(2) - np.exp(0.5)) / 3), 4))
\end{verbatim}
Keluaran sebenarnya:
\begin{verbatim}
nilai eigen A : [1. 1. 4.]
  sisa Av-lv  : 8e-16
  sisa Av-lv  : 1.1e-15
  sisa Av-lv  : 3e-16
A^10 lewat diagonalisasi :
[[349526 349525 349525]
 [349525 349526 349525]
 [349525 349525 349526]]
A^10 lewat pangkat biasa :
[[349526 349525 349525]
 [349525 349526 349525]
 [349525 349525 349526]]
selisih maks : 1.0477378964424133e-09
cond V       : 1.5227
iterasi pangkat pada A2:
  k=0 sudut=   0.00 rayleigh=4.0000
  k=1 sudut=  14.04 rayleigh=4.6471
  k=2 sudut=  21.25 rayleigh=4.8820
  k=3 sudut=  24.39 rayleigh=4.9578
  k=4 sudut=  25.69 rayleigh=4.9840
  k=5 sudut=  26.21 rayleigh=4.9938
sudut batas  : 26.57
eigen B      : [1. 1. 2.]
rank(B - I)  : 2
B^10         :
[[  -19    40     0]
 [  -10    21     0]
 [ 1023 -2046  1024]]
jumlah kolom : [1. 1. 1.]
eigen M      : [0.4 0.5 1. ]
x setelah 10 : [0.5672 0.2661 0.1666]
x setelah 40 : [0.5667 0.2667 0.1667]
eigen C      : [1. 2.]
eigen C+1e-6 : [0.999 2.001]
cond V(C)    : 2000.0
expm(0.5A)   : 1.9134
rumus        : 1.9134
\end{verbatim}
\end{lab}

Baris pertama keluaran itu menyebut nilai eigen dalam urutan menaik, $1,1,4$. Tiga baris
berikutnya adalah sisa $A\mathbf{v}-\lambda\mathbf{v}$ untuk ketiga vektor eigen yang dikembalikan
pustaka. Sisa berorde $10^{-16}$ berarti persamaannya terpenuhi sampai batas presisi mesin; tidak ada sisa
yang persis nol karena vektor eigen dinormalkan dan dihitung secara iteratif. Vektor eigen pustaka
juga bukan $(1,1,1)$ dan $(1,-1,0)$, melainkan kelipatan ternormalkan: arahnya sama, panjangnya
berbeda, dan diagonalisasi hanya peduli arah.

Dua blok $A^{10}$ memberi angka bulat yang identik, dengan selisih terbesar $1.05\times10^{-9}$.
Selisih itu bukan kesalahan aljabar: entri sebesar $3.5\times10^{5}$ dihitung lewat $V$, $V^{-1}$,
dan pangkat bilangan real, sehingga galat pembulatan berorde $10^{-16}$ ikut membesar.

Nilai \texttt{cond V} sebesar $1.5227$ dekat dengan satu, dan itu kabar baik: $A$ simetris,
sehingga vektor eigennya dapat dipilih saling tegak lurus dan $P$ mendekati ortogonal.
Diagonalisasi untuk matriks semacam itu jauh lebih stabil.

Blok iterasi pangkat mengulang hitungan pada Gambar~\ref{fig:iterasi-pangkat}: sudutnya naik dari
$0$ ke $26.21$ derajat, tetap di bawah batas $26.57$, sedangkan nisbah Rayleigh naik dari $4.0000$
ke $4.9938$ menuju $5$.

Blok $B$ menunjukkan sisi lain: nilai eigennya $1,1,2$ dan \texttt{rank(B - I)} sama dengan $2$,
jadi nullity-nya $1$. Multiplisitas geometrik untuk $\lambda=1$ hanya satu, sedangkan aljabarnya
dua; matriksnya defektif. Blok $B^{10}$ memperlihatkan akibatnya: entri $-19$ dan $1024$
berdampingan, yang pertama tumbuh linear dan yang kedua eksponensial.

Blok $M$ memeriksa bahwa jumlah setiap kolomnya satu, sehingga $\mathbf{x}_{k}$ tetap stokastik.
Nilai eigennya $0.4$, $0.5$, dan $1$; setelah empat puluh langkah angkanya berhenti bergerak pada
enam angka, sesuai laju $(1/2)^{k}$ yang diturunkan dengan tangan.

Blok terakhir memperingatkan satu hal: matriks $C$ punya nilai eigen $1$ dan $2$, dan bila satu entri
nol di bawah diagonal diganti $10^{-6}$, nilai eigennya bergeser menjadi $0.999$ dan $2.001$;
perpindahan $10^{-3}$, seribu kali lebih besar daripada gangguannya. Angka itu dapat dihitung
persis: akar $\lambda^{2}-3\lambda+1.999$ adalah $1.5\mp\sqrt{2510}/100$, yaitu $0.999000998$ dan
$2.000999002$. Pergeseran itu sebanding dengan $10^{-6}\cdot1000/(2-1)$: gangguan dikali entri besar
dibagi celah kedua nilai eigen. Bilangan kondisi $2000$ menangkap pembesaran yang sama. Nilai eigen
tetap kontinu terhadap entri, tetapi tidak harus kontinu secara Lipschitz dengan konstanta satu.

\begin{tantangan}
Carilah satu matriks $2\times2$ yang kedua nilai eigennya real dan berbeda, tetapi kedua vektor
eigennya tidak saling tegak lurus. Tunjukkan pilihan kita benar dengan menghitung hasil kali
titiknya, lalu jelaskan mengapa matriks simetris tidak dapat dipakai (jawabannya menunggu Bab 10).

Petunjuk 1: dua vektor tegak lurus tepat ketika hasil kali titiknya nol. Mulailah dari matriks
segitiga atas yang sederhana, lalu hitung kedua vektor eigennya tanpa membulatkan.

Petunjuk 2: matriks $\begin{pmatrix}2&1\\0&3\end{pmatrix}$ sudah cukup. Vektor eigen untuk $3$
adalah $(1,1)$ dan untuk $2$ adalah $(1,0)$; hasil kali titiknya $1$, bukan nol. Bandingkan dengan
matriks simetris $A$ di lab, yang memberi $\texttt{cond V}\approx1.5$.

Kalau mentok, kembali ke bagian Matriks serupa dan invarian, lalu ke paragraf bilangan kondisi di
Lab komputasi.
\end{tantangan}

\begin{untukdosen}
Urutan mengajar yang berhasil dipakai: nilai eigen lebih dulu sebagai arah yang dipertahankan,
baru polinomial karakteristik sebagai alat hitung. Bila polinomial datang lebih dulu, mahasiswa
melihat akar polinomial dan lupa bahwa yang dicari arah. Setelah ruang eigen, masuk multiplisitas
geometrik lewat nullity, baru syarat diagonalisasi. Contoh guntingan sebaiknya diberikan sebelum
teorema syarat, supaya teorema menjawab pertanyaan yang sudah terasa.

Dua miskonsepsi khas. Pertama, mahasiswa menganggap nilai eigen berulang selalu berarti tidak
dapat didiagonalkan. Matriks $(2,1,1;1,2,1;1,1,2)$ di bagian contoh membantahnya dengan
multiplisitas geometrik dua. Kedua, mereka menuliskan vektor eigen sebagai baris pada $P$, lalu
mendapat $P^{-1}AP$ yang bukan diagonal dan menyalahkan rumusnya. Aturan yang perlu diulang:
vektor eigen adalah kolom, dan urutan kolom $P$ harus sama dengan urutan diagonal $D$.

Aktivitas kelas: bagi kelas menjadi dua kelompok untuk matriks $(2,1,1;1,2,1;1,1,2)$. Kelompok
pertama menghitung $A^{10}$ lewat diagonalisasi, kelompok kedua lewat pangkat berulang atau
perkalian berulang cepat. Bandingkan hasilnya di papan, lalu minta kelas menjelaskan mengapa kedua
jalan harus bertemu. Kegaduhan biasanya muncul pada penyebut $3$ di $P^{-1}$.

Tanda kesiapan lanjut: mahasiswa dapat mengatakan, tanpa menghitung, mengapa matriks guntingan
tidak dapat didiagonalkan, dan dapat menebak bentuk $S^{k}$ untuk guntingan $2\times2$. Bila
mereka menyebut munculnya faktor $k$, mereka siap masuk Bab 10.
\end{untukdosen}

\section{Ringkasan}
\tingkat{dasar}

\begin{ringkasan}
\begin{itemize}
  \item $\lambda$ nilai eigen $A$ bila ada $\mathbf{v}\neq\mathbf{0}$ dengan
        $A\mathbf{v}=\lambda\mathbf{v}$. Nilai eigen adalah akar $p_{A}(\lambda)=\det(\lambda I-A)$,
        dan $E_{\lambda}=\operatorname{null}(A-\lambda I)$ adalah subruangnya.
  \item Derajat $p_{A}$ adalah $n$, jadi ada paling banyak $n$ nilai eigen berbeda. Atas
        $\mathbb{C}$ selalu ada $n$ bila dihitung dengan multiplisitas; atas $\mathbb{R}$ tidak,
        seperti pada matriks putar.
  \item $\operatorname{tr}(A)=\sum_{i}\lambda_{i}$ dan $\det(A)=\prod_{i}\lambda_{i}$, keduanya
        dihitung dengan multiplisitas. Ini pemeriksaan murah untuk setiap hitungan tangan.
  \item Untuk setiap nilai eigen, $1\le\dim E_{\lambda}\le$ multiplisitas aljabar. Vektor eigen
        dengan nilai eigen berbeda selalu bebas linear.
  \item $A$ dapat didiagonalkan tepat ketika ada $n$ vektor eigen bebas, dan tepat ketika
        multiplisitas aljabar sama dengan geometrik untuk setiap nilai eigen. Bila tidak, $A$
        disebut defektif, dan $A^{k}$ pada umumnya memuat faktor $k$, bukan hanya $\lambda^{k}$.
  \item Matriks serupa mempunyai polinomial karakteristik, nilai eigen, jejak, dan determinan
        yang sama; kebalikannya tidak berlaku. Vektor eigen ikut berubah bersama basis.
\end{itemize}
\end{ringkasan}

\section{Latihan}
\tingkat{menengah}

\begin{latihan}
\begin{enumerate}
  \item Tentukan nilai eigen dan vektor eigen
        $A_{1}=\begin{pmatrix}3&1\\ 2&2\end{pmatrix}$, lalu susun $P$ dan $D$ serta periksa
        $P^{-1}A_{1}P=D$ dengan mengalikan balik.

  \item Hitung $G^{10}$ untuk $G=\begin{pmatrix}4&2\\ 1&3\end{pmatrix}$ lewat diagonalisasi,
        lalu periksa hasilnya dengan fungsi \texttt{matrix\_power} dari NumPy. Tuliskan pula
        bentuk umum $G^{k}$.

  \item Buktikan bahwa bila $\lambda$ nilai eigen $A$ maka $\lambda^{2}$ nilai eigen $A^{2}$,
        dan secara umum $\lambda^{k}$ nilai eigen $A^{k}$ untuk setiap bilangan bulat $k\ge1$.
        Jelaskan mengapa kebalikannya gagal: berikan $A$ dengan $A^{2}=-I$, lalu bandingkan nilai
        eigen $A^{2}$ dengan nilai eigen real $A$.
        (Petunjuk: pakai matriks putar di awal bab; ia tidak punya nilai eigen real.)

  \item Buktikan bahwa $A$ invertible jika dan hanya jika $0$ bukan nilai eigen $A$, dan bahwa
        dalam hal itu nilai eigen $A^{-1}$ adalah $1/\lambda$ untuk setiap nilai eigen $\lambda$
        dari $A$.
        (Petunjuk: mulai dari $A\mathbf{v}=\lambda\mathbf{v}$, kalikan dengan $A^{-1}$ dari kiri,
        dan pastikan $\lambda\neq0$ supaya pembagiannya sah.)

  \item Nilai klaim ini: \emph{dua matriks persegi berukuran sama yang mempunyai polinomial
        karakteristik sama pasti serupa.} Tentukan benar, salah, atau belum cukup bukti, dan
        sertakan alasan atau contoh penyangkal.

  \item Buktikan bahwa $A$ dan $A^{\top}$ mempunyai polinomial karakteristik yang sama, dengan
        memakai $\det(\lambda I-A)=\det((\lambda I-A)^{\top})$ dan sifat determinan pada Bab 8.
        Jelaskan mengapa ini berarti transpos tidak mengubah nilai eigen, tetapi biasanya
        mengubah vektor eigen.

  \item Matriks $K=\begin{pmatrix}2&1&0\\ 0&2&0\\ 0&0&1\end{pmatrix}$ mempunyai nilai eigen
        $2,2,1$. Tentukan pangkat $K-2I$, lalu putuskan apakah $K$ dapat didiagonalkan. Jelaskan
        jawabanmu dengan syarat multiplisitas.

  \item Rantai Markov dua keadaan mempunyai matriks
        $T=\begin{pmatrix}0.9&0.4\\ 0.1&0.6\end{pmatrix}$. Tentukan nilai eigen $T$, titik
        tetapnya, dan banyak langkah yang diperlukan agar galat awal menyusut di bawah
        $10^{-6}$. Periksa jawabanmu dengan Python.
        (Petunjuk: setelah nilai eigen diketahui, galat menyusut seperti $|\lambda_{2}|^{k}$;
        selesaikan $|\lambda_{2}|^{k}<10^{-6}$ dengan logaritma.)

  \item Buktikan: bila $A$ dapat didiagonalkan dengan $A=PDP^{-1}$, maka
        $\operatorname{tr}(A^{k})=\sum_{i}\lambda_{i}^{k}$ untuk setiap $k\ge1$, dengan memakai
        $\operatorname{tr}(XY)=\operatorname{tr}(YX)$. Untuk $B$ yang defektif tadi, hitung
        $\operatorname{tr}(B^{10})$ dari matriks yang tercetak, lalu bandingkan dengan
        $1^{10}+1^{10}+2^{10}$. Jelaskan mengapa kesamaan itu tetap berlaku walaupun $B$
        defektif, dan di titik mana bukti tadi memakai diagonalisasi.
        (Petunjuk: jejak tidak berubah oleh perubahan basis; pakai bentuk $B^{k}$ yang sudah
        diturunkan.)
\end{enumerate}
\end{latihan}

\begin{untukpeneliti}
Yang benar-benar terbukti di bab ini adalah pernyataan aljabar atas medan apa pun dan berdimensi
hingga: kesetaraan nilai eigen dengan akar polinomial karakteristik, kebebasan vektor eigen untuk
nilai eigen berbeda, ketaksamaan $1\le\dim E_{\lambda}\le$ multiplisitas aljabar, dan kesetaraan
diagonalisasi dengan dua syarat yang dipakai di sini. Rujukan bakunya Axler, \emph{Linear Algebra
Done Right} edisi ke-4, bagian 5A dan 5D (5.55, empat syarat setara, dan 5.54); versi matriks ada
di Anton, \emph{Elementary Linear Algebra} edisi ke-11, bagian 5.1 dan 5.2.

Yang lebih sering dipopulerkan daripada dibuktikan ada dua. Pertama, klaim bahwa nilai eigen
``mengukur kepentingan'', dasar penjelasan PageRank: vektor eigen utama matriks stokastik memberi
bobot tunak hanya bila rantainya tak tereduksi dan aperiodik, dan itu isi teorema Perron-Frobenius,
yang tidak dibuktikan di sini; bukti lengkapnya di Strang, \emph{Introduction to Linear Algebra}
edisi ke-6, Apendiks 8. Kedua, klaim bahwa polinomial karakteristik adalah cara praktis menghitung
nilai eigen: untuk $n\ge5$ tidak ada rumus akar umum, dan pustaka memakai algoritma QR pada
matriks, bukan pemfaktoran polinomial. Bab ini juga berhenti pada kasus dapat didiagonalkan; bentuk
Jordan dan polinomial minimal tidak dibahas, dan ada di Hoffman dan Kunze, \emph{Linear Algebra}
edisi ke-2, bagian 6.3 dan 7.3.

Pertanyaan terbuka: seberapa jauh nilai eigen bergeser bila entri diganggu, yang jawaban
kuantitatifnya menuntut teorema Bauer-Fike atau pseudospektrum; dan apakah kedefektifan itu sendiri
stabil secara numerik. Batas keberlakuan: dimensi hingga, medan real atau kompleks, tanpa klaim
tentang stabilitas algoritma.
\end{untukpeneliti}

\peta{Bab ini mengikuti Anton, \emph{Elementary Linear Algebra} edisi ke-11, Bab 5: bagian 5.1
(Eigenvalues and Eigenvectors), 5.2 (Diagonalization), 5.3 (Complex Vector Spaces), 5.4
(Application: Differential Equations), dan 5.5 (Application: Dynamical Systems and Markov
Chains). Padanan yang lebih abstrak ada di Axler, \emph{Linear Algebra Done Right} edisi ke-4,
Bab 5: 5A (nilai eigen dan ruang eigen), 5B (keberadaan nilai eigen atas ruang vektor kompleks),
dan 5D (operator yang dapat didiagonalkan, termasuk 5.55). Untuk rantai Markov dan teorema
Perron-Frobenius, lihat Strang, \emph{Introduction to Linear Algebra} edisi ke-6, bagian 6.1,
6.2, 6.5, dan Apendiks 8. Bentuk kanonik untuk matriks defektif ada di Hoffman dan Kunze,
\emph{Linear Algebra} edisi ke-2, Bab 6 (Elementary Canonical Forms), bagian 6.3 untuk polinomial
minimal, dan Bab 7 (The Rational and Jordan Forms), bagian 7.3 untuk bentuk Jordan.}
